%% toc: P2 10 10.1,10.2
\section{Equivalence congruence, replacement, and the equivalence exercise}
\label{sec:p02-w02-p02-003-exercise-link}
\index{replacement!and constructor congruence}
\index{logical equivalence!congruence}

\Cref{ex:prop-metatheory-09} asks the reader to prove that logical
equivalence of \(\mathsf{CPL}\) formulas is an equivalence relation and a
congruence for every primitive connective.  The statement it asks for is now
part of a stated result, \cref{thm:p02-replacement}, whose complete proof is
Annex~\ref{annex:p2.prop.replacement}.  This section connects the exercise,
its feedback entry in \cref{app:exercise-feedback}, and the two results of
\cref{sec:p02-propositional-substitution} with their annex proofs.  It also
proves that, for any equivalence relation on formulas, the congruence asked
for in the exercise is exactly what is needed for replacement of selected
subformula occurrences.

\paragraph{Inputs and conventions.}
Formulas are the strings of \(\mathsf{Fm}_{\mathsf{CPL}}\) generated from
\(\mathsf{Prop}=\{p_n:n\in\mathbb N\}\) by the primitive constructors of
\cref{sec:prop-syntax-semantics-01}: the nullary connective \(\bottom\), the
unary connective \(\neg\), and the binary connectives \(\land,\lor,\to\).
Defined connectives such as \(\leftrightarrow\) are expanded before any
operation below is applied.  Valuations, their recursive extensions
\(\widehat v\), and logical equivalence \(\equiv\) are those of
\cref{sec:prop-syntax-semantics-02}: \(\alpha\equiv\beta\) means that every
valuation gives \(\alpha\) and \(\beta\) the same truth value.  Positions,
chosen occurrence sets \(O\), and the replacement \(\varphi[O\leftarrow\chi]\)
are those of \cref{def:p02-selected-replacement}; uniform substitution
\(\varphi[\upsilon]\) is that of \cref{def:p02-uniform-substitution}.

\begin{hbconvention}{def:p02-w02-p02-003-exercise-link}{Where the equivalence exercise is answered}
\Cref{ex:prop-metatheory-09} has two parts.
\begin{enumerate}
  \item Logical equivalence is reflexive, symmetric and transitive on
        \(\mathsf{Fm}_{\mathsf{CPL}}\).
  \item Logical equivalence is a congruence for every primitive connective.
        Since the alphabet lists \(\bottom\) as the nullary connective, this
        covers \(\bottom\equiv\bottom\) as well as the clauses for
        \(\neg\), \(\land\), \(\lor\) and \(\to\).
\end{enumerate}
Both parts are the first two assertions of \cref{thm:p02-replacement}, and
Annex~\ref{annex:p2.prop.replacement} proves them in full.  The theorem is
therefore a complete answer to the exercise.

The feedback entry for the exercise, headed \texttt{ex:prop-metatheory-09} in
the section on classical propositional logic of
\cref{app:exercise-feedback}, is a solution sketch.  Its steps correspond to
the annex proof as follows.
\begin{itemize}
  \item Reflexivity, symmetry and transitivity of \(\equiv\) come from the
        same three properties of equality of truth values, applied at each
        valuation.  This is the first paragraph of the annex proof.
  \item For congruence, the sketch fixes the inputs' truth values and uses
        the fact that each connective is a fixed truth function.  The annex
        proof fixes an arbitrary valuation \(v\), shows equal outputs at
        \(v\) from equal inputs at \(v\), and then uses that \(v\) was
        arbitrary.  It checks \(\neg\), \(\land\), \(\lor\) and \(\to\) one at
        a time.
  \item The sketch names the negation case and the binary cases.  The
        nullary case \(\bottom\equiv\bottom\), which part~2 also requires,
        is a separate line of the annex proof: every valuation assigns
        \(\mathbf F\) to \(\bottom\).
\end{itemize}
A written solution is complete when it supplies each of these items.

\Cref{thm:p02-replacement} goes beyond the exercise: if \(\psi\equiv\chi\),
then \(\varphi\equiv\varphi[O\leftarrow\chi]\) for every chosen set \(O\) of
occurrences of \(\psi\) in \(\varphi\).
\Cref{thm:p02-w02-p02-003-congruence-replacement} below shows that this
clause follows from the two parts of the exercise alone.  The companion
result for substituting formulas for variables is
\cref{thm:p02-uniform-substitution}, proved in
Annex~\ref{annex:p2.prop.substitution}.  Its last clause,
\(\varphi\equiv\psi\Rightarrow\varphi[\upsilon]\equiv\psi[\upsilon]\), says
that \(\equiv\) is also preserved when one uniform substitution is applied to
both sides.  The exercise does not ask for that clause.
\end{hbconvention}

\begin{hbtheorem}{thm:p02-w02-p02-003-congruence-replacement}{proved=annex:p02-w02-p02-003-replace}{Constructor congruence is equivalent to selected-occurrence replacement}
Let \(E\) be an equivalence relation on \(\mathsf{Fm}_{\mathsf{CPL}}\).  The
following are equivalent.
\begin{enumerate}
  \item[(a)] \(E\) is a congruence for every primitive connective:
        \(\bottom\mathrel{E}\bottom\); \(\alpha\mathrel{E}\alpha'\) implies
        \(\neg\alpha\mathrel{E}\neg\alpha'\); and whenever
        \(\alpha\mathrel{E}\alpha'\) and \(\beta\mathrel{E}\beta'\),
        \((\alpha\circ\beta)\mathrel{E}(\alpha'\circ\beta')\) for each
        \(\circ\in\{\land,\lor,\to\}\).
  \item[(b)] \(E\) admits selected-occurrence replacement: for all formulas
        \(\varphi,\psi,\chi\) with \(\psi\mathrel{E}\chi\) and every set
        \(O\subseteq\{a\in\operatorname{Pos}(\varphi):\varphi_a=\psi\}\),
        \(\varphi\mathrel{E}\varphi[O\leftarrow\chi]\).
\end{enumerate}
\end{hbtheorem}

For \(E={\equiv}\), part~(a) is the content of the exercise.  The direction
from (a) to (b) therefore turns any complete solution of the exercise into a
proof of the replacement clause of \cref{thm:p02-replacement}.  The direction
from (b) to (a) shows that nothing weaker than the exercise's congruence
suffices.  The proof uses reflexivity and transitivity of \(E\) and never
uses symmetry.

\begin{hbexample}{def:p02-w02-p02-003-worked-replacement}{One replacement checked through the exercise's clauses}
Let \(P,Q,R\) be distinct propositional variables, and put
\[
 \varphi=(\neg(P\to Q)\lor R),\qquad
 \psi=(P\to Q),\qquad
 \chi=(\neg P\lor Q).
\]
The positions of \(\varphi\) are \(\epsilon,1,11,111,112,2\), and
\(\varphi_{11}=\psi\) is the only occurrence of \(\psi\).  Take
\(O=\{11\}\).  The root is unselected, with \(O_1=\{1\}\) and
\(O_2=\varnothing\).  In \(\neg(P\to Q)\) the root is again unselected, now
with \(O_1=\{\epsilon\}\), so that subformula becomes \(\neg\chi\).  The
leaf \(R\) with the empty set is unchanged.  Hence
\[
 \varphi[O\leftarrow\chi]=(\neg(\neg P\lor Q)\lor R).
\]
The truth table of \cref{sec:prop-syntax-semantics-02} shows that the
columns for \(P\to Q\) and \(\neg P\lor Q\) agree on all four rows.  By
\cref{lem:prop-coincidence}, values of other variables do not matter, so
\(\psi\equiv\chi\).  The proof of (a)\(\Rightarrow\)(b) then runs through
three clauses of the exercise.  Negation congruence gives
\(\neg\psi\equiv\neg\chi\).  Reflexivity gives \(R\equiv R\).  Disjunction
congruence combines them into
\(\varphi\equiv\varphi[O\leftarrow\chi]\).  No other clause is used.

By \cref{thm:p02-uniform-substitution}, the equivalence
\(\psi\equiv\chi\) also survives any uniform substitution.  For example, with
\(\upsilon(P)=R\), \(\upsilon(Q)=\neg R\), and every other variable fixed,
\((R\to\neg R)\equiv(\neg R\lor\neg R)\).
\end{hbexample}

\begin{hbwarning}{def:p02-w02-p02-003-congruence-limits}{Two ways the congruence argument fails}
\emph{Agreement at one valuation is not equivalence.}  The feedback sketch
assumes that corresponding inputs have equal truth values.  This must be
read at one fixed valuation at a time, and the conclusion holds at that
valuation only.  An equivalence follows only after every valuation has been
treated.  Agreement at a single valuation is strictly weaker.  For example, a
valuation with \(v(P)=v(Q)=\mathbf T\) gives \(P\) and \(Q\) the same value,
but \(P\not\equiv Q\): a valuation \(w\) with \(w(P)=\mathbf T\) and
\(w(Q)=\mathbf F\) separates them, and it also separates \(\neg P\) from
\(\neg Q\).

\emph{An equivalence relation need not be a congruence.}  Call formulas
\(\alpha\) and \(\beta\) equisatisfiable when \(\{\alpha\}\) is satisfiable
if and only if \(\{\beta\}\) is.  This is an equivalence relation on
\(\mathsf{Fm}_{\mathsf{CPL}}\).  The formulas \(P\) and \(\neg P\) are
equisatisfiable, since both are satisfiable, and \(P\) is equisatisfiable
with itself.  Yet \((P\land P)\) is satisfiable and \((P\land\neg P)\) is
not, so conjunction congruence fails.  By the direction (b)\(\Rightarrow\)(a)
of \cref{thm:p02-w02-p02-003-congruence-replacement}, replacement must
therefore fail for this relation.
\Cref{def:p02-substitution-replacement-limits} gives an explicit
instance.
\end{hbwarning}

\paragraph{Relationships.}
\Cref{thm:p02-replacement,thm:p02-uniform-substitution} are the stated
results.  \Cref{ex:prop-metatheory-09} remains an exercise, and its feedback
entry remains a solution sketch.  The definitions of position, chosen
occurrence set and replacement are those of
\cref{def:p02-selected-replacement}, and
\cref{def:p02-substitution-replacement-example} contrasts replacement with
uniform substitution on one formula.

Kleene proves a deductive replacement theorem in the propositional calculus:
from the hypothesis that two formulas are equivalent in his formal sense, the
equivalence of a formula with the result of replacing one specified
occurrence of the first by the second is deducible.  His proof is an
induction on the depth of that occurrence, and its steps are lemmas, one for
each operator and argument place, that carry such an equivalence through
that operator \cite[\S\,26, Theorem~6 and lemmas *26--*30, p.~116]{Kleene1952}.
He also notes that the theorem contains those lemmas as its cases of depth~1,
and that a set of occurrences is replaced by iterating single replacements
\cite[\S\,26, p.~117]{Kleene1952}.  His induction corresponds to the
direction (a)\(\Rightarrow\)(b) of
\cref{thm:p02-w02-p02-003-congruence-replacement}, and the containment of
the lemmas corresponds to (b)\(\Rightarrow\)(a).  The
relation here is instead the semantic relation \(\equiv\), or an arbitrary
equivalence relation \(E\), and a chosen set of occurrences is replaced
simultaneously in one recursion.
